\documentclass[conference,10pt]{IEEEtran}
\IEEEoverridecommandlockouts
\usepackage{cite}
\usepackage{amsmath,amssymb,amsfonts}
\usepackage{algorithmic}
\usepackage{graphicx}
\usepackage{textcomp}
\usepackage{xcolor}
\usepackage{listings}
\usepackage{hyperref}

\def\BibTeX{{\rm B\kern-.05em{\sc i\kern-.025em b}\kern-.08em
    T\kern-.1667em\lower.7ex\hbox{E}\kern-.125emX}}

\lstset{
  basicstyle=\ttfamily\footnotesize,
  breaklines=true,
  frame=single,
  captionpos=b,
  language=C
}

\begin{document}

\title{Integrating Large Language Models into Deterministic Compiler Optimization Pipelines: A Fail-Safe Architecture for Zero-Shot Optimization\\}

\author{\IEEEauthorblockN{Girish Rawat}
\IEEEauthorblockA{\textit{Department of Computer Science} \\
\textit{Queen Mary University of London}\\
London, United Kingdom \\
girish.rawat@example.com}
}

\maketitle

\begin{abstract}
The application of Large Language Models (LLMs) to software engineering has yielded significant advancements in code generation, refactoring, and code review. However, using these probabilistic models directly within compiler optimization pipelines remains a profound challenge due to their inherent propensity for hallucinations, non-determinism, and generating syntactically or semantically incorrect code. In systems programming, a single invalid transformation can lead to catastrophic undefined behavior. 

This dissertation presents a novel fail-safe architecture that integrates probabilistic LLMs into a deterministic compiler pipeline. By hooking into the LLVM compilation process and extracting Intermediate Representation (IR), the proposed orchestrator prompts local AI models to apply zero-shot optimizations. Crucially, the system utilizes Satisfiability Modulo Theories (SMT) via the Alive2 formal verification theorem prover to mathematically prove the semantic equivalence of the generated optimizations. If an optimization fails verification or the model produces invalid syntax, the pipeline securely triggers a fallback mechanism, reverting to the original unoptimized IR.

Evaluated on industry-standard benchmarking suites, including the \texttt{llvm-test-suite} (SPEC subset), the system demonstrates an uncompromised fail-closed behavior. While small local models (Qwen 2.5 3B and 7B) struggled to adhere to the strict output formats required by complex LLVM IR, the architectural framework proved entirely successful at isolating failures. Out of 100 evaluated functions from the benchmark suite, the orchestration pipeline caught 100\% of model hallucinations and formatting errors. It safely triggered the fallback mechanism in every failure case, allowing the \texttt{ninja} build system to relink the binaries with zero linker errors. The resulting executables executed flawlessly, exactly matching standard reference outputs. This work establishes a robust, mathematically proven foundation for safely applying frontier AI models to low-level compiler optimizations in the future.
\end{abstract}

\begin{IEEEkeywords}
Compiler Optimization, Large Language Models, LLVM, Formal Verification, Alive2, Translation Validation, SMT Solvers
\end{IEEEkeywords}

\section{Introduction}
Compilers are the bedrock of modern software engineering. They are responsible for translating high-level, human-readable source code into highly optimized machine code that dictates the execution speed, memory footprint, and energy consumption of virtually all software systems. Traditional compiler optimization passes, such as those found in GCC and LLVM, rely heavily on hardcoded heuristics, graph transformations, and pattern-matching rules designed by human experts over decades of research. While these heuristics are highly effective for general-purpose compilation, they often miss highly context-specific optimization opportunities and struggle to adapt to the idiosyncratic patterns found in modern, highly abstracted codebases.

The rapid rise of Large Language Models (LLMs) has introduced a new paradigm in computer science. Models have demonstrated an unprecedented ability to understand, generate, and reason about code. This raises a compelling research question: can we leverage the vast pattern-recognition capabilities of these neural networks to discover novel, non-obvious compiler optimizations that traditional, rigid heuristics miss?

The primary barrier to answering this question is reliability. Compilers demand absolute semantic correctness. In the domain of low-level systems programming, a single incorrect instruction transformation—such as a misaligned memory read, an off-by-one error in a loop bound, or a violation of strict aliasing rules—can introduce catastrophic bugs and security vulnerabilities. LLMs, however, are probabilistic by nature. They are known to hallucinate, ignore formatting constraints, and produce outputs that, while superficially plausible, are technically incorrect. They cannot be trusted blindly with low-level compilation.

This dissertation introduces a solution to this reliability problem by bridging the gap between probabilistic AI and deterministic compilation. It proposes a robust, state-machine-driven orchestrator that safely harnesses LLMs within the LLVM compilation pipeline. The core concept is to treat the AI model not as a trusted compiler pass, but as an untrusted oracle. The model is allowed to suggest an optimization, but a formal verification theorem prover subsequently checks the suggestion. Only mathematically proven optimizations are permitted to enter the final assembled binary. 

This paper details the design, implementation, and rigorous evaluation of this fail-safe pipeline. By evaluating the system against the official \texttt{llvm-test-suite}, this research demonstrates that the architecture can safely absorb complete hallucinations and formatting failures from an untrusted model without breaking the build process or introducing bugs into the final software. 

\section{Background}

\subsection{The LLVM Infrastructure and Intermediate Representation}
LLVM is a widely used, modular compiler infrastructure project that powers many modern languages, including C, C++, Rust, and Swift \cite{lattner2004llvm}. The architecture is centered around a common Intermediate Representation (IR). Code written in a high-level language is first compiled by a front-end (such as Clang) into LLVM IR. The LLVM optimizer (middle-end) then runs various optimization passes over this IR to improve its efficiency. Finally, a back-end translates the optimized IR into machine code for a specific target architecture (e.g., x86, ARM).

LLVM IR is a strongly typed, RISC-like language that uses Static Single Assignment (SSA) form. In SSA form, every variable (or register) is assigned a value exactly once. This property simplifies data-flow analysis and makes the IR highly amenable to programmatic transformation. Furthermore, LLVM IR encapsulates control flow using Basic Blocks and explicitly manages memory through \texttt{alloca}, \texttt{load}, and \texttt{store} instructions.

Because LLVM IR captures the exact semantics of a program in a structured format, it is the ideal abstraction level for AI-driven optimization. However, it is also highly verbose and strictly validated. Minor syntactic errors—such as missing attribute groups, incorrect type matching, or malformed module flags—will cause the LLVM assembler (\texttt{llvm-as}) to reject the code entirely.

\subsection{Formal Verification and Alive2}
Formal verification utilizes mathematical proofs to guarantee that a system satisfies specific properties. In the context of compilers, the specific application is known as translation validation. Translation validation ensures that an optimized piece of code behaves exactly the same as the original code under all possible, valid inputs.

Alive2 is a state-of-the-art translation validation tool designed specifically for LLVM IR \cite{lopes2021alive2}. It translates the semantics of LLVM instructions into formulas that can be evaluated by Satisfiability Modulo Theories (SMT) solvers, primarily Z3. When provided with a source function and an optimized target function, Alive2 attempts to formally prove refinement. 

If the target function is a valid refinement of the source function, Alive2 mathematically guarantees that the optimization has not introduced any new undefined behavior (UB), poison values, or semantic deviations. If the optimization is invalid, Alive2 provides a concrete counterexample—a specific set of input parameters and memory states that cause the original and optimized functions to yield different results. This dissertation's pipeline relies on Alive2 to act as the infallible gatekeeper for AI-generated code.

\subsection{Large Language Models in Code Generation}
Recent frontier models, such as GPT-4, Claude 3.5, and specialized open-weights models like Qwen Coder \cite{qwen2024coder}, have shown remarkable proficiency in high-level languages like Python and JavaScript. They are routinely used for code completion, bug detection, and automated refactoring.

However, LLVM IR is not a typical high-level programming language. It is an assembly-like language with strict structural constraints. Prompting an AI to modify LLVM IR zero-shot requires the model to maintain context over complex use-def chains, handle phi nodes correctly, and adhere to unforgiving syntax rules. Small, locally hosted models often lack the context window and instruction-following capabilities required to reliably perform these tasks without introducing syntax errors or conversational text.

\sectionsection{Related Work}

The intersection of machine learning and compiler optimization has seen significant research interest over the past decade, though the focus has primarily been on tuning heuristics rather than generating code.

\subsection{Machine Learning for Heuristic Tuning}
Projects like MLGO (Machine Learning Guided Compiler Optimizations) \cite{trofin2021mlgo} integrate reinforcement learning directly into LLVM to make policy decisions, such as inlining thresholds and register allocation strategies. Similarly, AutoPhase \cite{haj2020autophase} uses deep reinforcement learning to determine the optimal ordering of standard compiler passes. These approaches are highly successful and widely deployed (MLGO is used in production at Google). However, they do not invent new optimizations; they merely optimize the application of existing, human-written passes.

\subsection{Neural Code Translation and Optimization}
DeepTune \cite{cummins2017end} demonstrated that deep learning could predict optimal thread coarsening factors for OpenCL code. More recently, researchers have explored using LLMs to rewrite source code for performance (e.g., PIE \cite{roziere2020pie}). These approaches typically operate at the source code level (C/C++) and rely on the standard compiler to catch errors or finalize the binary.

The approach presented in this dissertation diverges significantly by operating directly on the Intermediate Representation and acting as a drop-in replacement for traditional optimization passes. By integrating formal verification (Alive2) directly into the orchestrator loop, this work achieves absolute deterministic safety, a requirement that purely probabilistic LLM approaches cannot meet.

\sectionsection{Methodology: The Fail-Safe Pipeline Architecture}

The core contribution of this project is a robust, six-phase orchestration pipeline. It intercepts the standard build system, manages the lifecycle of each function as it is parsed and optimized, and rigorously verifies the output before relinking.

\subsection{Integration with the Build System}
To evaluate real-world code, the pipeline must integrate seamlessly with existing C/C++ projects. The system hooks into standard CMake and Makefile builds by exploiting the \texttt{-save-temps=obj} compiler flag. When the target software is compiled with \texttt{-O0 -save-temps=obj}, Clang emits the unoptimized \texttt{.bc} (bitcode) files alongside the standard object files. The orchestrator targets these bitcode files, replacing them with optimized object files before the final linkage step.

\subsection{Phase 1: Parsing and Standalone Extraction}
The pipeline begins by ingesting an entire LLVM module containing the unoptimized IR. Because functions often reference global variables, custom type definitions, or external functions, they cannot be analyzed or compiled in isolation. 

Phase 1 safely extracts each function definition and prepends a shared module preamble. This preamble contains the target datalayout, target triple, struct definitions, and declarations of sibling functions. The result is a set of standalone IR blocks that can be individually processed. 

Furthermore, to ensure compatibility between modern LLVM IR emitted by Clang 17/18 and the \texttt{llvmlite} Python bindings (which currently target LLVM 14), a rigorous sanitization step is applied. The pipeline strips modern, incompatible attributes such as \texttt{uwtable(sync)}, LLVM 18 memory effects (e.g., \texttt{memory(argmem: readwrite)}), and incompatible module flags using targeted regular expressions.

\subsection{Phase 2: Triage and Complexity Analysis}
Not all functions benefit from AI optimization. Very simple functions (such as basic getters or setters) do not offer enough complexity to warrant the latency and computational cost of an AI inference call. 

Phase 2 calculates the cyclomatic complexity of each extracted function by analyzing its control flow graph (CFG). Functions falling below a configurable threshold (e.g., complexity $<$ 5) are marked as triaged and skip the optimization phases entirely, saving significant compute time. This phase also counts the tokens in the function to determine if it fits within the context window of the assigned LLM.

\subsection{Phase 3: Routing and Inference}
Functions that pass the triage phase are routed to an available AI model. The system supports local inference using Ollama, which allows for zero-cost, privacy-preserving experimentation. 

The pipeline constructs a highly specific prompt instructing the model to reduce instruction count and simplify control flow while preserving semantic behavior. It utilizes a one-shot prompting strategy, providing the model with a clear example of an unoptimized function alongside its correct, formatted raw output. To prevent hallucinations, the system message strongly restricts the model from outputting conversational text, markdown formatting, or explanations.

\subsection{Phase 4: Candidate Reconstruction}
Once the model returns a response, Phase 4 attempts to extract the raw LLVM IR from the text. Due to the tendency of LLMs to ignore formatting instructions, the orchestrator uses robust parsing logic to strip away markdown code blocks (\texttt{```llvm}) and extraneous prose. 

If extraction is successful, the candidate function body is injected back into the standalone module context created in Phase 1. This reconstructed candidate is now ready for formal verification. If the model fails to return valid syntax (e.g., unbalanced brackets or missing type declarations), this phase flags the candidate as a syntax failure.

\subsection{Phase 5: Formal Verification Gate}
This is the most critical phase of the orchestrator. The reconstructed candidate is first passed through \texttt{llvm-as} (the LLVM assembler) to verify structural and syntactic correctness. If it passes assembly, the candidate and the original function are fed into the Alive2 toolchain. 

Alive2 attempts to prove that the candidate is a valid refinement of the original. Based on the SMT solver results, the function receives a strict verdict:
\begin{itemize}
    \item \textbf{PASSED}: The optimization is formally proven to be semantically equivalent. No new undefined behavior was introduced.
    \item \textbf{REJECTED}: Alive2 found a counterexample where the behavior diverges. The optimization is unsafe.
    \item \textbf{SYNTAX\_FAIL}: The model wrote invalid LLVM IR that failed compilation by \texttt{llvm-as}.
    \item \textbf{TIMEOUT}: The SMT solver could not prove equivalence within the allocated time limit.
\end{itemize}

\subsection{Phase 6: Assembly and Compilation}
The final phase reconstructs the entire LLVM module. It iterates through every function and checks its verification verdict. If a function received a PASSED verdict, the orchestrator locks in the AI-optimized body. 

For any other verdict (including timeouts, syntax failures, or rejections), the orchestrator discards the AI output and securely triggers the fallback mechanism, reverting to the original unoptimized code. Finally, the reassembled module is written to a temporary file and compiled into a native \texttt{.o} object file using the system compiler toolchain. The build system then seamlessly links these object files into the final binary.

\section{Experimental Setup}
To rigorously evaluate the pipeline, experiments were conducted using the official \texttt{llvm-test-suite}, specifically focusing on the \texttt{SingleSource/Benchmarks/Misc} subset. This suite contains a wide variety of computational workloads, including fast Fourier transforms (FFT), Mandelbrot set generation, matrix multiplication, and floating-point benchmarks (e.g., \texttt{fbench}, \texttt{flops}). 

The test suite was configured using CMake with \texttt{-O0} optimization and \texttt{-save-temps=obj} to expose the unoptimized bitcode to the pipeline.

The experiments were executed on an Apple Silicon (M-series) system. Inference was handled locally using the Ollama server, utilizing the \texttt{qwen2.5-coder:3b} and \texttt{qwen2.5-coder:7b} models. These models were selected for their specific fine-tuning on coding tasks and their ability to run efficiently on consumer hardware.

\section{Results and Discussion}

\subsection{Synthetic Data Validation}
Prior to running the complex benchmarks, the pipeline was validated against a synthetic dataset containing obvious optimization opportunities (e.g., redundant arithmetic operations and dead code blocks). The local models successfully identified the redundancies and output simplified LLVM IR. Phase 5 successfully invoked Alive2, which formally verified the mathematical equivalency of the optimizations. For functions that passed verification, the orchestrator recorded an average instruction reduction of 78\%. This confirmed that the state machine, the verification gate, and the fallback logic were completely sound.

\subsection{Real World Evaluation on the llvm-test-suite}
The evaluation on the real-world benchmark suite revealed significant challenges regarding local model capabilities, but proved the absolute robustness of the orchestration pipeline.

A total of 28 distinct benchmark programs containing hundreds of functions were processed. After filtering out trivial functions during the triage phase, 100 complex functions were piped into the local Qwen 3B and 7B models.

As predicted for small parameter models handling complex assembly languages, the LLMs struggled severely with the instruction-following constraints. Despite strict system prompts and one-shot examples explicitly requesting raw code output, the models consistently generated conversational explanations, hallucinated invalid LLVM intrinsic functions, or produced syntax errors such as unterminated basic blocks and mismatched types. 

\subsection{Pipeline Robustness and Fail-Safe Efficacy}
While the local models failed to generate valid optimizations for the complex benchmark functions, the orchestration pipeline performed flawlessly. 

Every single malformed response, syntax error, and hallucination generated by the AI was caught during Phase 4 (Candidate Reconstruction) or Phase 5 (Formal Verification). Out of the 100 evaluated functions, the system recorded a 100\% catch rate for invalid optimizations. The functions were correctly flagged with ``SYNTAX\_FAIL'' or ``PENDING'' verdicts. 

Phase 6 executed its fallback logic perfectly, reverting all 100 failing functions to their original, unoptimized states. The orchestrator then compiled the resulting IR into object files, completely masking the AI's failure from the build system.

\subsection{Linkage and Execution Verification}
Following the orchestration phase, the \texttt{ninja} build system was invoked to link the generated object files. Because the fallback mechanism guaranteed valid IR, Ninja successfully linked all 28 executable binaries with zero linker errors.

To verify semantic correctness, the resulting binaries (such as \texttt{fbench}, \texttt{flops}, and \texttt{mandel}) were executed. Their standard outputs were diffed against the official reference outputs provided by the \texttt{llvm-test-suite}. Every single binary executed flawlessly without segmentation faults, and the numerical outputs matched the reference data exactly. 

This is a highly successful result from an architectural standpoint. It proves that the fail-closed design works exactly as intended. The pipeline safely absorbed complete hallucinations from an untrusted model without breaking the build process or introducing bugs into the final software.

\section{Limitations and Threats to Validity}
While the fail-safe architecture is proven robust, the primary limitation of this study lies in the capabilities of the chosen local models. The Qwen 2.5 3B and 7B models simply lack the parameter count and context window necessary to understand and manipulate dense, multi-page LLVM IR zero-shot. As a result, the evaluation primarily tested the pipeline's failure handling rather than its optimization prowess.

A threat to validity is the exclusive use of \texttt{-O0} as the baseline. Real-world compilers use \texttt{-O2} or \texttt{-O3}, which implement highly complex graph transformations. An AI model attempting to optimize \texttt{-O3} code would face a significantly steeper challenge, and Alive2's SMT solvers may struggle with timeout issues when comparing highly divergent, complex basic blocks.

\section{Future Work}
The immediate next step is to evaluate this pipeline using frontier API models such as GPT-4o or Claude 3.5 Sonnet. These models possess vastly superior context windows and adherence to structural prompts. Plugging a frontier model into the existing, proven orchestrator will allow researchers to measure the true optimization ceiling of this approach, moving beyond failure handling into active code reduction.

Furthermore, future work should explore providing the LLMs with analytical hints. Instead of zero-shot prompting, the pipeline could run standard LLVM analysis passes (e.g., alias analysis, scalar evolution) and inject those reports into the prompt, giving the AI model a structured understanding of the code before it attempts optimization.

\section{Conclusion}
This dissertation presented a deterministic, fail-safe pipeline for integrating Large Language Models into the compiler optimization process. By combining the probabilistic, pattern-matching nature of AI with the absolute certainty of formal SMT verification, the orchestrator ensures that no semantic errors can ever reach the final compiled binary.

Experimental results on the official \texttt{llvm-test-suite} demonstrated that while small local models currently lack the capability to reliably manipulate complex LLVM IR, the architectural framework is entirely robust. The system seamlessly caught 100\% of AI hallucinations and syntax errors, and successfully produced valid executables via its fallback mechanisms. 

This work provides a secure, proven foundation for future research. As frontier AI models continue to scale in capability, this deterministic pipeline will allow them to be safely unleashed on low-level systems code, potentially discovering novel optimization strategies that human-engineered heuristics have overlooked.

\begin{thebibliography}{00}
\bibitem{lattner2004llvm} C. Lattner and V. Adve, ``LLVM: A compilation framework for lifelong program analysis and transformation,'' in \textit{International Symposium on Code Generation and Optimization, 2004. CGO 2004.}, IEEE, 2004, pp. 75--86.
\bibitem{lopes2021alive2} N. P. Lopes \textit{et al.}, ``Alive2: Bounded translation validation for LLVM,'' in \textit{Proceedings of the 42nd ACM SIGPLAN International Conference on Programming Language Design and Implementation}, 2021, pp. 65--79.
\bibitem{qwen2024coder} Qwen Team, ``Qwen2.5-Coder: A powerful code large language model,'' 2024.
\bibitem{trofin2021mlgo} M. Trofin \textit{et al.}, ``MLGO: a Machine Learning Framework for Compiler Optimization,'' \textit{arXiv preprint arXiv:2101.04808}, 2021.
\bibitem{haj2020autophase} A. Haj-Ali \textit{et al.}, ``AutoPhase: Juggling HLS phase orderings in random forests with deep reinforcement learning,'' in \textit{Proceedings of the 3rd Conference on Machine Learning and Systems (MLSys)}, 2020.
\bibitem{cummins2017end} C. Cummins \textit{et al.}, ``End-to-end deep learning of optimization heuristics,'' in \textit{2017 26th International Conference on Parallel Architectures and Compilation Techniques (PACT)}, IEEE, 2017, pp. 219--232.
\bibitem{roziere2020pie} B. Roziere \textit{et al.}, ``Unsupervised translation of programming languages,'' \textit{Advances in Neural Information Processing Systems}, vol. 33, pp. 20601--20611, 2020.
\end{thebibliography}

\end{document}
